\documentclass[12pt,letterpaper]{article}
%\setlength{\oddsidemargin}{.0in}
%\setlength{\evensidemargin}{.0in}
%\setlength{\textwidth}{6.5in}
%\setlength{\topmargin}{-.3in}
%\setlength{\headsep}{.20in}
%\setlength{\textheight}{9.in}
\usepackage[margin=1in]{geometry}
\usepackage[leqno]{amsmath}
\usepackage{amsfonts}
\usepackage{amssymb}
\usepackage{amsthm, amscd}
\usepackage{amssymb}
\usepackage[all]{xy}
\usepackage{graphicx}
\usepackage{enumerate}


%Here are some user-defined notations
\newcommand{\RR}{\mathbf R}  %bold R
\newcommand{\CC}{\mathbf C}  %bold C
\newcommand{\ZZ}{\mathbf Z}   %bold Z
\newcommand{\QQ}{\mathbf Q}   %bold Q
\newcommand{\rr}{\mathbb R}     %blackboard bold R
\newcommand{\cc}{\mathbb C}    %blackboard bold R
\newcommand{\zz}{\mathbb Z}    %blackboard bold R
\newcommand{\qq}{\mathbb Q}   %blackboard bold Q
\newcommand{\ZZn}[1]{\ZZ/{#1}\ZZ}
\newcommand{\zzn}[1]{\zz/{#1}\zz}
\newcommand{\calM}{\mathcal M}  %calligraphic M
\newcommand{\sm}{\setminus} 
\newcommand{\bfa}{\mathbf a}
\newcommand{\bfb}{\mathbf b}
\newcommand{\bfc}{\mathbf c}


%improving spacing in tables (space above and below characters in a row)
\newcommand{\tfix}{\rule{0pt}{2.6ex}}
\newcommand{\bfix}{\rule[-1.2ex]{0pt}{0pt}}


%Here are commands with variable inputs 
\newcommand{\intf}[1]{\int_a^b{#1}\,dx}
\newcommand{\intfb}[3]{\int_{#1}^{#2}{#3}\,dx}
\newcommand{\ip}[2]{\left\langle {#1}, {#2}\right\rangle}
\newcommand{\pln}[1]{$\sm${\tt #1}}
\newcommand{\bgn}[1]{$\tt {\sm}begin\{#1\}$}
\newcommand{\nd}[1]{$\tt {\sm}end\{#1\}$}
\newcommand{\marginalfootnote}[1]{%
        \footnote{#1}
        \marginpar[\hfill{\sf\thefootnote}]{{\sf\thefootnote}}}
\newcommand{\edit}[1]{\marginalfootnote{#1}}


%Here are some user-defined operators
\newcommand{\Tr}{\operatorname {Tr}}
\newcommand{\nul}{\operatorname {null}}
\newcommand{\ran}{\operatorname {ran}}
\renewcommand{\span}{\operatorname {span}}
\newcommand{\GL}{\operatorname {GL}}
\newcommand{\SL}{\operatorname {SL}}
\newcommand{\Prob}{\operatorname {Prob}}
\newcommand{\re}{\operatorname {Re}}   %new definition of "real part" operator
\newcommand{\im}{\operatorname {Im}}   %new definition of "imaginary part" operator


%These commands deal with theorem-like environments (i.e., italic)\nul
\theoremstyle{plain}
\newtheorem{theorem}{Theorem}[section]
\newtheorem{proposition}{Proposition}[section]
\newtheorem{corollary}[theorem]{Corollary}
\newtheorem{lemma}[theorem]{Lemma}
\newtheorem{conjecture}[theorem]{Conjecture}

%These deal with definition-like environments (i.e., non-italic)
\theoremstyle{definition}
\newtheorem{definition}[theorem]{Definition}
\newtheorem{example}[theorem]{Example}
\newtheorem{remark}[theorem]{Remark}

%This numbers equations by section
\numberwithin{equation}{section}


%This is for hypertext references
\usepackage{color}
\usepackage{hyperref}


%spacing
\usepackage{setspace}




\begin{document}
\onehalfspacing
\tableofcontents
\addtocontents{toc}{~\hfill\textbf{Page}\par}

\section{Basic Notions}
\begin{definition} A binary operation on a set $S$ is a function $f:S\times S\rightarrow S$. 
\end{definition}
\begin{definition} A \emph{field} is a set $\mathbb{F}$ together with two binary operations $+$, and $\cdot$ called addition and multiplication (respectively) such that 
\ \\
\begin{enumerate}[1.]
\setlength{\itemsep}{5pt}
\item For all $a,b,c\in \mathbb{F}$ we have 
\[a+(b+c)=(a+b)+c\]
and 
\[a\cdot(b\cdot c)=(a\cdot b)\cdot c.\]
\item For all $a,b\in \mathbb{F}$ we have 
\[a+b=b+a\]
and 
\[a\cdot b=b\cdot a.\]
\item There exists an element $0\in \mathbb{F}$, called an additive identity,  such that for all $a\in \mathbb{F}$ we have $a+0=a$. 
\item There exists an element $1\in \mathbb{F}$, called a multiplicative identity, such that for all $a\in \mathbb{F}$ we have $a\cdot 1=a$.
\item For all $a\in \mathbb{F}$ there exists an element $b\in \mathbb{F}$, called an additive inverse, such that $a+b=0.$ 
\item For all $a\in \mathbb{F}$ such that $a\neq 0$ there exists an element $c\in \mathbb{F}$, called a multiplicative inverse, such that $a\cdot c=1$.
\item  For all $a,b,c\in \mathbb{F}$ 
\[a\cdot (b+c)=(a\cdot b)+ (a\cdot c)\]

\end{enumerate}
\end{definition}
\ \\
{\bf \noindent Note:} Fields have a unique additive and multiplicative identity denoted 0 and 1 respectively. Moreover, when the additive and multiplicative inverses exist they are unique.\\

\begin{example}
 All of the following examples are with their standard operations. 
\begin{enumerate}[1.]
\setlength{\itemsep}{5pt}
\item $\mathbb{Q}$ (rational numbers)
\item $\mathbb{R}$ (real numbers)
\item $\mathbb{C}$ (complex numbers)
\item $\mathbb{Z}/p\mathbb{Z}$ for $p$ prime (Integers modulo p)
\end{enumerate}
\end{example}

{\bf\noindent Non example: }$\mathbb{Z}$ is not a field, it lacks multiplicative inverses. 

\begin{definition} A \emph{vector space} $V$ over a field $\mathbb{F}$ is a set $V$ with two operations called \emph{vector addition} and \emph{scalar multiplication} where vector addition is a function $+:V\times V\rightarrow V$ and scalar multiplication is a function $\cdot: \mathbb{F}\times V\rightarrow V$ such that 
\\
\begin{enumerate}[1.] 
\setlength{\itemsep}{5pt}
\item For all $u,v\in V$ we have 
\[u+v=v+u\]
\item For all $u,v,w\in V$ and for all $a,b\in\mathbb{F}$ we have 
\[(u+v)+w=u+(v+w)\]
and 
\[(ab)\cdot v=a\cdot(b\cdot v)\]
\item There exists a vector $0\in V$, called an additive identity, such that for all $v\in V$ we have
\[v+0=v\]

\item For all $v\in V$ we have a vector $w\in V$, called an additive inverse, such that 
\[v+w=0\]
\item For all $v\in V$ we have 
\[1\cdot v=v\]
\item For all $a,b \in \mathbb{F}$ and for all $u,v\in V$ we have
\[a\cdot(u+v)=a\cdot u+a\cdot v\] 
\end{enumerate}
\end{definition}

\begin{example} All of the following examples are with their standard operations. 
\begin{enumerate}[1.]
\setlength{\itemsep}{5pt}
\item $\mathbb{F}^n=\left\{\begin{bmatrix}a_1\\ \vdots \\ a_n\end{bmatrix}: a_i\in \mathbb{F}\right\}$ where $\mathbb{F}$ is a field. 
\item Polynomials with coefficients in a field $\mathbb{F}$.
\item Polynomials ( with coefficients in a field $\mathbb{F}$) of degree $\leq n$
\item Continuous functions $f:X\rightarrow Y$, $C(X,Y)$, where $X$ and $Y$ are fields. 
\item Functions from a field $X$ into a field $Y$. 
\item $\mathbb{F}^\infty=\{(a_1, a_2, a_3, \ldots): a_i\in \mathbb{F}\}$. 
\end{enumerate}
\end{example}

\begin{proposition}Every vector space $V$ has a unique additive identity. The unique additive identity is denoted $0$. 
\end{proposition}

\begin{proposition}Every element $v\in V$ has a unique additive inverse. For all $v\in V$ its unique additive inverse is denoted $-v$. 
\end{proposition}

\begin{proposition}For all $v\in V$ we have $0\cdot v=0$. 
\end{proposition}

\begin{proposition}For all $a\in \mathbb{F}$ and $0\in V$ we have $a\cdot 0=0$.
\end{proposition}

\begin{proposition}For every $v\in V$ we have $(-1)\cdot v=-v$
\end{proposition}

\section{Basis for a Vector Space}

\begin{definition} A \emph{linear combination} of a list of vectors $v_1,\ldots, v_m$ in $V$ is a vector of the form 
\[a_1v_1+\ldots+a_mv_m\]
where $a_1,\ldots, a_m\in \mathbb{F}$. 
\end{definition}

\begin{definition} The set of all linear combinations of a list of vectors $v_1, \ldots, v_m$ in $V$ is called the \emph{span} of $v_1, \ldots,v_m$ denoted by 
$\span\{v_1,\ldots,v_m\}$.
\[\span\{v_1,\ldots,v_m\}=\{a_1v_1+\ldots+a_mv_m\mid a_i\in \mathbb{F}\}\]
\end{definition}

\begin{definition} If $V$ is a vector space and $V=\span\{v_1,\ldots,v_m\}$ then we say that $v_1, \ldots, v_m$ span $V$. 
\end{definition}

\begin{definition} We say that a vectors space is \emph{finite dimensional} if there exists a finite list of vectors $v_1,\ldots, v_m$ such that 
\[\span\{v_1,\ldots v_m\}=V\] 
Otherwise we say that $V$ is \emph{infinite dimensional}.
\end{definition}


\begin{definition} A list of vectors $v_1,\ldots, v_m$ in $V$ is called \emph{linearly independent} if the only choice of $a_1, \ldots, a_m\in \mathbb{F}$ such that 
\[a_1v_1+\ldots+a_mv_m=0\] is $a_1=a_2=\ldots =a_m$. A list is called \emph{linearly dependent} if it is not linearly independent.
\end{definition}

\begin{lemma} Suppose that $v_1,\ldots, v_m$ is a linearly dependent list in $V$. There exists a $j\in\{1,\ldots, m\}$ such that \\
\begin{enumerate}[1)]
\item $v_j\in \span\{v_1,\ldots v_{j-1}\}$
\item $\span\{v_1,\ldots, v_{j-1},v_j,v_{j+1},\ldots v_m\}=\span\{v_1,\ldots, v_{j-1},v_{j+1},\ldots v_m\}$
\end{enumerate}
\end{lemma}

\begin{proposition} In a finite dimensional vector space the length of every linearly independent list of vectors is less than or equal to the length of every spanning list of vectors. 
\end{proposition}


\begin{definition} A basis for a vector space $V$ is a list of vectors $\{v_1, \ldots, v_n\}$ such that 
\begin{enumerate}[1.]
\item $\{v_1, \ldots, v_n\}$ is linearly independent
\item $\span\{v_1, \ldots, v_n\}=V$. 
\end{enumerate}
\end{definition}
\begin{proposition} A list of vectors $\{v_1, \ldots, v_n\}$ in $V$ is a basis for $V$ if and only if every vector $v\in V$ can be written \emph{uniquely} in the form
\[v=a_1v_1+\ldots+a_nv_n\]
for some $a_1, \ldots, a_n\in \mathbb{F}$. 
\end{proposition}
\begin{proposition} Every spanning list of vectors in $V$ can be reduced down to a basis. 
\end{proposition}

\begin{proposition} Every linearly independent list of vectors in $V$ can be extended to a basis.  
\end{proposition}

\begin{proposition} Any two basis of a finite dimensional vector space $V$ have the same length. 
\end{proposition}

\begin{definition} The dimension of a finite dimensional vector space is the length of any basis of the vector space. Denoted $\dim(V)$. 
\end{definition}

\begin{proposition} Suppose $V$ is finite dimensional. Every list of linearly independent vectors whose length is equal to the dimension of $V$ is a basis. 
\end{proposition}

\begin{proposition} Suppose $V$ is finite dimensional. Every spanning list vectors whose length is equal to the dimension of $V$ is a basis. 
\end{proposition}

\section{Subspaces}
\begin{definition} A subspace of a vector space $V$ is a subset $H$ such that $H$ is a vector space under the same binary relations and field as $V$. 
\end{definition}

\begin{proposition}[Subspace Test] A subset $H$ is a subspace of $V$ if and only if 
\begin{enumerate}[1.] 
\item $0\in H$. 
\item For all $u,v\in H$ we have $u+v\in H$
\item For all $u\in H$ and $a\in \mathbb{F}$ we have $au\in H$.
\end{enumerate}
\end{proposition}

\begin{proposition} If $U$ is a subspace of a finite dimensional vector space $V$ then $\dim(U)\leq\dim(V)$. Moreover, $\dim(U)=\dim(V)$ if and only if $V=U$. 
\end{proposition}

\subsection{Direct Sums}
\begin{definition} Suppose $U_1,\ldots, U_m$ are subsets of $V$. The \emph{sum} of $U_1, \ldots, U_m$ denoted $U_1+\ldots+U_m$ is the set of all possible sums i.e., 
\[U_1+\ldots+U_m=\{u_1+\ldots+u_m\ |\ u_i\in U_i, i=1,\ldots, m\}\]
\end{definition}

\begin{proposition} If $U_1, \ldots, U_m$ are subspaces then so is $U_1+\ldots+U_m$. 
\end{proposition}

\begin{definition} Suppose $U_1,\ldots, U_m$ are subspaces of $V$. The sum $U_1+\ldots+U_m$ is a \emph{direct sum} if each element of $U_1+\ldots+U_m$ can be written in only one way as a sum $u_1+\ldots +u_m$ where $u_i\in U_i$, $i=1,\ldots,m$. The direct sum is denoted $U_1\oplus \ldots \oplus U_m$. 
\end{definition}

\begin{proposition}  $U_1+\ldots+U_m$ is a direct sum if and only if the only way to write $0$ as a sum is by taking each $u_i$ where $i=1, \ldots, m$ to be $0$. 
\end{proposition}

\begin{proposition}  The sum of two subspaces $U$ and $W$ is a direct sum if and only if\\ $U\cap W=\{0\}$. 
\end{proposition}

\begin{proposition} If $V$ is a finite dimensional vector space and $U$ is a subspace of $V$ then there exists a $W$ which is a subspace of $V$ such that $V=U\oplus W$. 
\end{proposition}

\subsection{Quotient Spaces} 

\begin{definition}Let $V$ be a vectors space and $U$ a subspace. For every $v\in V$ define 
\[v+U=\{v+u\ \mid\  u\in U\}\]
and 
\[V/U=\{v+U\ \mid\  v\in V\}\]
\end{definition}
\begin{proposition}Let $V$ be a vectorspace, $U$ a subspace and $v,w\in V$.  The following are equivalent. 
\begin{enumerate}[(a)]
\item $v-w\in U$
\item $v+U=w+U$
\item $(v+U)\cap (w+U)\neq \emptyset$
\end{enumerate}
\end{proposition}

\begin{proposition} Let $V$ be a vectorspace, $U$ a subspace, $\lambda\in \mathbb{F}$, and $v,w\in V$The set $V/U$ is a vector space with the following operations: 
\[(v+U)+(w+U)=(v+w)+U\]
\[\lambda(v+U)=(\lambda v)+U\]
\end{proposition}


\begin{proposition} Let $V$ be a finite dimensional vector space and $U$ be a subspace. 
\[\dim(V/U)=\dim(V)-\dim(U)\]
\end{proposition}

\section{Linear Maps}
\begin{definition} A linear map from $V(\mathbb{F})$ to $W(\mathbb{F})$ is a function $T:V\rightarrow W$ such that 
\[T(\lambda x+y)=\lambda T(x)+T(y)\]
for every $x,y\in V$ and $\lambda\in \mathbb{F}$. Denote the set of all linear maps from $V$ to $W$ as $\mathcal{L}(V,W)$. 
\end{definition}
\begin{proposition} Suppose $\{v_1,\ldots, v_n\}$ is a basis of $V$ and $w_1,\ldots,w_n\in W$. There is a unique linear map $T:V\rightarrow W$ such that 
\[T(v_j)=w_j\ \quad \ j=1,\ldots, n.\]
\end{proposition}

\begin{proposition} If $T:V\rightarrow W$ is linear then 
\[T(0_V)=0_W.\]
\end{proposition}

\begin{definition} Suppose $S, T\in \mathcal{L}(V,W)$ and $\lambda \in \mathbb{F}$ define 
\[(S+T)(v)=S(v)+T(v) \text{ for all }v\in V\]
and
\[(\lambda T)(v)=\lambda T(v) \text{ for all }v\in V.\]
\end{definition}
\begin{proposition} $\mathcal{L}(V,W)$ is a vector space with the above operations.
\end{proposition}
\subsection{One to One, Onto, Invertibility, and Isomorphisms}
\begin{definition} For $T\in \mathcal{L}(V,W)$ define the \emph{null space} of $T$ (or kernel of T) to be 
\[\nul(T)=\{v\in V\ \mid \  Tv=0\}\subseteq V\]
and define the \emph{range} of $T$ to be 
\[\ran(T)=\{Tv\ \mid \  v\in V\}\subseteq W.\]
\end{definition}

\begin{theorem} For $T\in \mathcal{L}(V,W)$ both $\nul(T)$ and $\ran(T)$ are subspaces. A linear transformation is injective if and only if 
\[\nul(T)=\{0\}.\]
A linear transformation is surjective if and only if $\ran(T)=W$. 
\end{theorem}
\begin{definition} The \emph{rank} of a linear transformation is the dimension of its range. The \emph{nullity} of a transformation is the dimension of its null space. 
\end{definition}

\begin{definition} A vector space isomorphism from $V$ onto $W$ is a bijective linear map $T:V\rightarrow W$. If there is a vector space isomorphism from $V$ onto $W$ we say $V$ is isomorphic to $W$ and write $V\cong W$. 
\end{definition}
\begin{theorem} Let $V$ and $W$ be two finite dimensional vector spaces. $V$ is isomorphic to $W$ if and only if $\dim(V)=\dim(W)$.
\end{theorem}
\begin{definition} A linear map $T:V\rightarrow W$ is \emph{invertible} if there exists an map $S:W\rightarrow V$ such that 
\[S\circ T=\text{Id}_{V}\]
\[T\circ S=\text{Id}_{W}\]
\end{definition}
\begin{theorem} A map $T:V\rightarrow W$ is invertible if and only if its bijective. 
\end{theorem}
\subsection{Fundamental Theorem of Linear Maps}
\begin{proposition} Suppose $T\in \mathcal{L}(V,W)$ and define 
\[\tilde{T}:V/\nul(T)\rightarrow W\] 
by
\[\tilde{T}(v+\nul(T))=Tv.\]
The following hold:
\begin{enumerate}[a)]
\item $\tilde{T}$ is linear
\item $\tilde{T}$ is injective
\item $\ran(\tilde{T})=\ran(T)$
\item $V/\nul(T)\cong \ran(T)$.
\end{enumerate}

\end{proposition}
\begin{theorem}[Fundamental Theorem of Linear Maps/ Rank-Nullity]
Suppose $V$ is a finite dimensional vector space and $T\in \mathcal{L}(V,W)$. We have $\ran(T)$ is finite dimensional and 
\[\dim(V)=\dim(\ran(T))+\dim(\nul(T)).\]
\end{theorem}

\begin{proposition} Suppose $V$ and $W$ are finite dimensional vector spaces. Let $T\in \mathcal{L}(V,W)$. 
\begin{enumerate}[a)]
\item If $\dim(V)<\dim(W)$ then $T$ is not surjective.
\item If $\dim(V)>\dim(W)$ then $T$ is not injective.
\end{enumerate}
\end{proposition}

\begin{definition} A linear map $T\in \mathcal{L}(V,V)$ is called an \emph{operator}. 
\end{definition}
\begin{theorem} Suppose $T$ is an operator over a vector space $V$. If $V$ is finite dimensional the following are equivalent:
\begin{enumerate}[a)]
\item $T$ is injective
\item $T$ is surjective
\item $T$ is bijective
\end{enumerate}
\end{theorem}

\subsection{The Matrix of a Linear Map and the Coordinate Transform}
\begin{definition} Let $V$ be an $n$-dimensional vector space over the field $\mathbb{F}$. Let $\beta=\{v_1, \ldots, v_n\}$ be an \emph{ordered} basis for $V$. The coordinate transform $\varphi_\beta:V\rightarrow \mathbb{F}^n$ is defined by 
\[v=a_1v_1+\ldots+a_nv_n\overset{\varphi_\beta}{\longmapsto} \begin{bmatrix}a_1\\ \vdots \\ a_n\end{bmatrix}\]
Denote the column vector $\varphi_\beta(v)$ by $[v]_\beta$.
\end{definition}
\begin{proposition} The coordinate transform from $V$ onto $\mathbb{F}^n$ is a vector space isomorphism.
\end{proposition}

\begin{definition} Suppose $T:V\rightarrow W$ is a linear maps with $\beta=\{v_1,\ldots, v_n\}$ and $\gamma=\{w_1, \ldots, w_m\}$ as ordered basis for $V$ and $W$ respectively. Define the matrix of $T$, denoted by $[T]_\beta^\gamma$, by the following
If 
\[T(v_k)=a_{1,k}w_1+\ldots +a_{m,k}w_k\]
then the $k$-th column of $[T]_\beta^\gamma$ is given by $\begin{bmatrix} a_{1,k}\\ \vdots \\ a_{m,k}\end{bmatrix}$. 
\end{definition}

\begin{theorem} The following diagram commutes\\
\[
\begin{CD}
V @>T>> W\\
@V\phi_{\beta}VV @AA\phi_\gamma^{-1} A\\
\mathbb{F}^n @>[T]_{\beta}^\gamma>> \mathbb{F}^m
\end{CD}
\]
More specifically, 
\[[T]_{\beta}^\gamma[v]_{\beta}=[T(v)]_{\gamma}.\]
\end{theorem}
\section{Determinants}
\begin{definition} Let $\hat{A}_{i,j}$ be the $(n-1)\times (n-1)$ matrix that results from $A$ by removing the $i$th row and $j$th column and let $a_{i,j}$ be the entry in the $i$th row and $j$th column.  Consider the set of $n\times n$ matrices over $\mathbb{F}$.  
Define \[\det([a])=a\]
and
\[\det(A)=\sum_{i=1}^n(-1)^{i+1}a_{i,1}\det(\hat{A}_{i,1})\]
\end{definition}
\subsection{Multilinear and Alternating}\label{Multilinear and Alternating}

\begin{definition} Let $V_1, \ldots V_n$ be vector spaces over a field $\mathbb{F}$. The product $V_1\times \ldots \times V_n$ is defined by 
\[V_1\times \ldots \times V_n=\{(v_1,\ldots v_n)\ \mid \ v_1\in V_1, \ldots v_n\in V_n\}\] 
\end{definition}
Of course, with the appropriate operations $V_1\times \ldots \times V_n$ is a vector space. 
\begin{proposition} $V_1\times \ldots \times V_n$ is a vector space over $\mathbb{F}$ with the following operations:
\[(v_1,\ldots, v_n)+(u_1, \ldots, u_n)=(v_1+u_1, \ldots, v_n+u_n)\]
\[c(v_1, \ldots, v_n)=(cv_1, \ldots, cv_n).\]
\end{proposition}
The proof of the above proposition is standard, we will omit it. 
\begin{definition} Let $V_1, \ldots ,V_n, W$ be vector spaces over a field $\mathbb{F}$. 
A map $\varphi:V_1\times\ldots \times V_n\rightarrow W$ is called \emph{multilinear} if for each fixed $i$ and fixed elements $v_j\in V_j$, $j\neq i$, the map
\[V_i\rightarrow W\quad \text{ defined by }\quad x\mapsto \varphi(v_1, \ldots, v_{i-1}, x, v_{i+1}, \ldots ,v_n)\]
is linear. If each $V_i=V$, $i=1,2,\ldots n$ then $\varphi$ is called a $n$-multilinear function on $V$. If $W$ is a field, then $\varphi$ is called a multilinear form on $V$. 

The function $\det:M_{n\times n}(\mathbb{F})\rightarrow \mathbb{F}$ is viewed as a multilinear map by viewing the columns of a matrix as column vectors and making the following identification. 
\[M_{n\times n}(\mathbb{F})\ni A=[v_1,\ldots, v_n]\mapsto (v_1, \ldots, v_n)\in \mathbb{F}^n\times \ldots \times \mathbb{F}^n.\]
\end{definition}

\begin{proposition}$\det:M_{n\times n}(\mathbb{F})\rightarrow \mathbb{F}$ is a multilinear function (viewing each matrix as a tuple of column vectors in $\mathbb{F}^n\times \ldots \times \mathbb{F}^n$).
\end{proposition}

\begin{definition}An $n$-multilinear function $\varphi$ on $V$ is called alternating if $\varphi$ is zero whenever two consecutive arguments are equal, i.e. if $v_i=v_{i+1}$ for some $i\in\{1,\ldots, n-1\}$, then $\varphi(v_1, \ldots ,v_n)=0$. 
\end{definition}


\begin{lemma}Let $B\in M_{n\times n}(\mathbb{F})$, where $n\geq 2$. If the column $j$ of $B$ equals 
\[e_k=\left[\begin{array}{c}0\\\vdots\\ 1\\\vdots\\ 0\end{array}\right] - \text{ k}\small{th} \text{ spot}\] for some $k$ $(1\leq k\leq n)$, then 
\[\det(B)=(-1)^{j+k}\det(\hat{B}_{k,j})\]
\end{lemma}

\begin{theorem} The determinant of a square matrix can be evaluated by cofactor expansion along any column or row, i.e. 
\[\det(A)=\sum_{i=1}^n (-1)^{i+j}a_{i,j} \det(\hat{A}_{i,j})\]
or 
\[\det(A)=\sum_{j=1}^n (-1)^{i+j}a_{i,j} \det(\hat{A}_{i,j})\]
\end{theorem}


\begin{proposition}
The determinant function $\det:M_{n\times n}(\mathbb{F})\rightarrow \mathbb{F}$ is an alternating function.
\end{proposition}

\begin{theorem}\label{unique}
The determinant $\det:M_{n\times n}(\mathbb{F})\rightarrow \mathbb{F}$ is the \emph{unique} multilinear alternating map taking the identity matrix to the mulitiplicative identity element in $\mathbb{F}$
\end{theorem}


\subsection{The Multiplicative Property}


\begin{definition}[Elementary Row Operations]
An elementary row operation is any one of the following operations performed on a matrix. 
\begin{enumerate}[$\bullet$]
\item Switching the position of two rows.
\item Multiplying the entries of a row by a scalar.
\item Replacing a row with its addition of a scalar multiple of another row. 
\end{enumerate}
\end{definition} 
\begin{lemma}\label{ele prod}
Let $A$ and $B$ be matrices such that $C=AB$ is defined. Suppose $e_1, \ldots e_n$ be a sequence of elementary row operations. Let $A^\prime$ be the matrix resulting from performing $e_1,\ldots, e_n$ on $A$ and $C^\prime$ be the matrix resulting from performing $e_1, \ldots, e_n$ on $C$. Then 
\[C^\prime=A^\prime B\]
\end{lemma}

\begin{lemma}\label{det ele}
 Suppose $e_1, \ldots e_n$ be a sequence of elementary row operations. Let $A^\prime$ be the matrix resulting from performing $e_1,\ldots, e_n$ on $A$. We have that 
 \[\alpha\det(A^\prime)= \det{A}\]
 for some $\alpha\in \mathbb{F}$ depending only on $e_1,\ldots, e_n$. 
\end{lemma}

\begin{definition}
Let $A$ be a matrix. Define its transpose, denoted $A^\top$, by
\[(A^\top)_{i,j}=A_{j,i}\]
\end{definition}

\begin{lemma}
Let $A$ and $B$ be matrices such that $AB$ is defined. 
\[(AB)^\top=B^\top A^\top\]
\end{lemma}

\begin{lemma}\label{det trans}
For any square matrix $A$ 
\[\det(A)=\det(A^\top)\]
\end{lemma}


\begin{definition} A lower triangular matrix is any matrix $L$ such that 
\[L_{i,j}=0 \quad\text{ for }\quad j>i\]
\end{definition}

\begin{definition} An upper triangular matrix is any matrix $U$ such that 
\[U_{i,j}=0 \quad\text{ for }\quad j<i\]
\end{definition}

\begin{lemma} The product of upper (lower) triangular matrices is an upper (lower) triangular matrix.
\end{lemma}

\begin{lemma}\label{det tri}
Let $A$ be either an upper triangular matrix or lower triangular matrix. We have that 
\[\det(A)=\prod a_{i,i}.\]
\end{lemma}


\begin{proposition}\label{multiplicativeid}
\[\det(AB)=\det(A)\det(B)\]
\end{proposition}

\subsection{Invertibility of a Matrix}
\begin{definition}A matrix $A\in M_{n\times n}(\mathbb{F})$ is said to be invertible if there exists a matrix $B\in M_{n\times n}(\mathbb{F})$ such that 
\[AB=I\]
\[BA=I\]
where $I$ is the $n\times n$ identity matrix. 
\end{definition}

\begin{definition}
Suppose $S$ is a set with a binary operation $\cdot :S\times S\rightarrow S$, then $S$ with $\cdot$ is a monoid if 
\begin{enumerate}[(a)]
\item For all $a, b$ and $c$ in $S$ 
\[(a\cdot b)\cdot c=a\cdot (b\cdot c)\]
\item There exists an identity element $1$ in $S$ such that for all $a\in S$
\[a\cdot 1=1\cdot a=a.\]
\end{enumerate}
\end{definition}

\begin{definition}
Let $A$ be a monoid and $1$ its identity element. We say $a\in A$ is invertible if and only if there exists a $b\in A$ such that
\[a\cdot b=b\cdot a=1\]
\end{definition}

\begin{definition}
A monoid homomorphism between monoids $A$ and $B$ is a function $\phi:A\rightarrow B$ such that 
\[\phi(a_1)\cdot \phi(a_2)=\phi(a_1\cdot a_2)\]
for all $a_1, a_2\in A$ 
and \[\phi(1_A)=1_B\] for the identity elements $1_A\in A$ and $1_B\in B$.
\end{definition}

\begin{theorem}
Let $A$ and $B$ be two monoids and $\phi:A\rightarrow B$ a monoid homomorphism. If $a\in A$ is invertible then $\phi(a)$ is invertible. 
\end{theorem}


\begin{theorem}\label{matinv}
Let $T\in \mathcal{L}(V,W)$ and let $\beta=\{v_1,\ldots v_n\}$, $\gamma=\{w_1, \ldots w_m\}$ be bases for $V$ and $W$ respectively. The transformation $T$ is invertible if and only if $[T]_\beta^\gamma$ is invertible.
\end{theorem}

\begin{definition}
An elementary matrix is any matrix obtained from performing a single row operation on the identity matrix. 
\end{definition}


\begin{proposition} Let $A$ be a matrix and let $B$ be the matrix obtained from $A$ by performing a row operation with corresponding elementary matrix $E$. We have 
\[EA=B\]
\end{proposition}


\begin{proposition}If $A$ is a matrix and $B$ is an upper triangular matrix obtained from $A$ via a finite sequence of row operations, then $A$ is invertible if and only if $B$ is invertible. Moreover, $B$ is invertible if and only if the entries on the diagonal are non-zero. 
\end{proposition}


\begin{theorem} A matrix $A$ is invertible if and only if $\det(A)\neq 0$. 
\end{theorem}

\subsection{Properties and Facts}
\begin{enumerate}[$\bullet$]
\item If $A$ is a triangular matrix then $\det(A)$ is the product of the entries on the main diagonal.
\item If a multiple of one row of $A$ is added to another row to produce $B$ then $\det(A)=\det(B)$.
\item If two rows of $A$ are interchanged to produce $B$ then $\det(A)=-\det(B)$. 
\item If one row of $A$ is multiplied by $k$ to produce $B$ then $\det(B)=k\cdot \det(A)$. 
\item A square matrix is invertible if and only if $\det(A)\neq 0$.
\item $\det(A^\top)=\det(A)$. 
\item $\det(AB)=\det(A)\det(B)$.
\item $\det(I)=1$.
\end{enumerate}
\section{Eigenvalues and Eigenvectors}
{\color{red}{\bf \noindent Note: }We will assume all vector spaces in this section are complex! (However, a lot of the following is definable for a vector spaces over any field)}

\begin{definition} Suppose $T\in \mathcal{L}(V)$. A subspace $U$ of $V$ is called \emph{invariant under} $T$ if $u\in U$ implies that $Tu\in U$. 
\end{definition}
\begin{definition}Suppose $T\in \mathcal{L}(V)$. A scalar $\lambda\in \mathbb{F}$ is an \emph{eigenvalue} of $T$ if there exists a $v\in V$ such that $v\neq 0$ and $Tv=\lambda v$. The vector $v$ is called an \emph{eigenvector}.
\end{definition}
\begin{proposition}$T\in \mathcal{L}(V)$ has a one dimensional subspace if and only if $T$ has an eigenvalue. 
\end{proposition}

\begin{theorem}Suppose $V$ is a finite dimensional vector space and $T\in \mathcal{L}(V)$. The following are equivalent:
\begin{enumerate}[(a)]
\item $\lambda\in \mathbb{F}$ is an eigenvalue.
\item $T-\lambda Id$ is not injective.
\item $T-\lambda Id$ is not surjective.
\item $T-\lambda Id$ is not invertible.
\end{enumerate}
\end{theorem}

\begin{proposition} Suppose  $T\in \mathcal{L}(V)$. If $\lambda_1, \ldots, \lambda_m$ are distinct eigenvalues of $T$ and $v_1, \ldots, v_m$ are corresponding eigenvectors then $\{v_1, \ldots, v_m\}$ is a linearly independent set. 
\end{proposition}
\begin{corollary} Suppose $T\in \mathcal{L}(V)$ and $\dim(V)=n$. Then $T$ has at most $n$ distinct eigenvalues. 
\end{corollary}

\begin{definition}Suppose  $T\in \mathcal{L}(V)$. Define for positive $m$
\[T^m=\underbrace{T\circ\ldots\circ T}_{m\text{ times}}\]
\[T^0=Id\]
and if $T$ is invertible
\[T^{-m}=\underbrace{T^{-1}\circ\ldots\circ T^{-1}}_{m\text{ times}}\]
\end{definition}

\begin{definition}Suppose  $T\in \mathcal{L}(V)$ and let $p(x)=a_nx^n+\ldots+a_1x+a_0$ be a polynomial over $\mathbb{F}$. Define
\[p(T)=a_nT^n+\ldots+a_1T+a_0Id\]

\end{definition}

\begin{proposition} Let $T\in \mathcal{L}(V)$. If $p(x)$ is a polynomial over $\mathbb{F}$ and $p(x)=s(x)q(x)$ then $p(T)=s(T)q(T)$. 
\end{proposition}


\begin{theorem} Every operator on a finite dimensional, non-zero, \emph{complex} vector space has an eigenvalue. 
\end{theorem}

\subsection{Upper Triangular Matrices}
\begin{proposition}Suppose $T\in \mathcal{L}(V)$ and $\beta=\{v_1, \ldots, v_n\}$ is a basis of $V$. The following statements are equivalent. 
\begin{enumerate}[(a)]
\item $[T]_\beta$ is upper triangular. 
\item $Tv_j\in \span\{v_1, \ldots, v_n\}$ for all $1\leq j\leq n$. 
\item $\span\{v_1, \ldots, v_j\}$ is invariant under $T$ for all $1\leq j\leq n$. 
\end{enumerate}
\end{proposition}

\begin{theorem}Suppose $T\in \mathcal{L}(V)$ where $V$ is finite dimensional. The operator $T$ has an upper triangular matrix with respect to some basis. 
\end{theorem}

\begin{proposition}Suppose $T\in \mathcal{L}(V)$ and $\beta=\{v_1, \ldots, v_n\}$ is a basis of $V$. If $[T]_\beta$ is an upper triangular matrix then the following statements hold:
\begin{enumerate}[(a)]
\item $T$ is invertible if and only if the diagonal entries of $[T]_\beta$ are non-zero. 
\item The eigenvalues of $T$ (and $[T]_\beta$) are the diagonal entries of $[T]_\beta$. 
\item $\det([T]_\beta)$ is the product of the eigenvalues of $T$. 
\item If $\gamma$ is another basis for $V$ then $\det([T]_\gamma)$ is the product of the eigenvalues of $T$. 
\end{enumerate}
\end{proposition}
\subsection{Diagonal Matrices}
\begin{proposition}Suppose $T\in \mathcal{L}(V)$ where $V$ is $n$-dimensional. If $T$ has $n$ distinct eigenvalues then $T$ has a diagonal matrix with respect to some basis of $V$. 
\end{proposition}

\begin{theorem} Suppose $T\in \mathcal{L}(V)$ where $V$ is finite dimensional. Let $\lambda_1, \ldots, \lambda_m$ be the distinct eigenvalues of $T$. The following are equivalent. 
\begin{enumerate}[(a)]
\item $T$ has a diagonal matrix with respect to some basis. 
\item $V$ has a basis of eigenvectors
\item There exists one dimensional subspaces $U_1, \ldots, U_n$ of $V$ where $U_i=\span\{v_i\}$ $1\leq i\leq n$ and each $U_i$ is invariant under $T$ such that $V=U_1\oplus \ldots \oplus U_n$. 
\item $V=\nul(T-\lambda_1Id)\oplus \ldots \oplus \nul(T-\lambda_m Id)$ 
\item $\dim(V)=\dim(\nul(T-\lambda_1Id))+ \ldots  + \dim(\nul(T-\lambda_m Id))$ 
\end{enumerate}
\end{theorem}

\section{Inner Products} 
{\color{red}{\bf \noindent Note: }We will assume all vector spaces in this section are either real or complex!}

\begin{definition}Let $V=\mathbb{C}^n$ (or $V=\mathbb{R}^n$). We define the \emph{dot product} of two vectors $\vec{x}=\left[x_1,\ldots ,x_n\right]^\top$ and $\vec{y}=\left[y_1,\ldots ,y_n\right]^\top$ as 
\[\vec{x}\cdot \vec{y}=x_1\bar{y}_1+\ldots+x_n\bar{y}_n\] where if $z=a+ib\in \mathbb{C}$ for $a,b\in \mathbb{R}$ then $\bar{z}=a-ib$. We call $\bar{z}$ the \emph{complex conjugate} of $z$. 
\end{definition}


\begin{definition}Let $V$ be a complex vector space (or a real vector space). An inner product on $V$ is a function 
\[\ip{\cdot}{\cdot}:V\times V\rightarrow \mathbb{C} \quad \quad(\text{or }\ip{\cdot}{\cdot}:V\times V\rightarrow \mathbb{R}) \]
such that 
\begin{enumerate}[(a)]
\item $\ip{v}{v}\geq 0$ for all $v\in V$. 
\item $\ip{v}{v}=0$ if and only if $v=0$. 
\item $\ip{\lambda v+u}{w}=\lambda\ip{v}{w}+\ip{u}{w}$ for all $\lambda \in \mathbb{C} \ (\text{or } \mathbb{R})$ and $u,v,w\in V$. 
\item $\ip{v}{w}=\overline{\ip{w}{v}}$ \quad \quad(or $\ip{v}{w}=\ip{w}{v}$).
\end{enumerate}
A vector space $V$ with an inner product defined on $V$ is called an \emph{inner product space}.
\end{definition}
{\bf \noindent Note: }Properties (3) and (4) say that the inner product is linear in the first variable and conjugate linear in the second variable. This is the common convention in mathematics. Physics uses the opposite convention on inner products, i.e. the inner product is conjugate linear in the first variable and linear in the second variable. This is mostly a matter of convenience for physicists as they use the ``bra-ket" notation for inner-products. \emph{{ In this course we will use the mathematics convention!}}

\begin{example}
The following are inner product spaces. 
\begin{enumerate}[(a)]
\item $\mathbb{C}^n$ (or $\mathbb{R}^n$) with the dot product. 
\item $\ell^2(\mathbb{C})=\left\{{\bf a}=(a_0, a_1, \ldots )\middle | \sum_{n=0}^\infty |a_i|^2<\infty \right\}$ with $\ip{{\bf a}}{{\bf b}}=\sum_{i=0}^\infty a_i\overline{b_i}$. 
\item $\mathcal{L}^2(\mathbb{R}, dx)=\left\{ f:\mathbb{R}\rightarrow \mathbb{C} \middle | \int_{-\infty}^\infty |f(x)|^2\, dx<\infty \right\}$ with $\ip{f}{g}=\int_{-\infty}^\infty f(x)\overline{g}(x)\, dx$. 
\end{enumerate}
\end{example}

\begin{definition} The \emph{norm} or \emph{length} of a vector $v$ in an inner product space $V$ is defined to be $\|v\|=\sqrt{\ip{v}{v}}$.
\end{definition}

\begin{definition} If $z=a+ib\in \mathbb{C}$ define the \emph{modulus} (or absolute value) of a complex number to be $|z|:=\sqrt{a^2+b^2}$.
\end{definition}

\begin{lemma} For $v\in V$ and $a\in \mathbb{F}$ we have $\|av\|=|a|\|v\|$. 
\end{lemma}

\subsection{Geometry}
\begin{definition} We say that two vectors $v$ and $u$ in an inner product space $V$ are \emph{perpendicular} or \emph{orthogonal} if $\ip{v}{u}=0$. We say two vectors $v$ and $u$ are \emph{parallel} if $u=av$ for some $a\in \mathbb{F}$. 
\end{definition}
\begin{theorem}[Pythagorean Theorem] 
If $u$ and $v$ are orthogonal vectors in an inner product space $V$ then \[\|u+v\|^2=\|u\|^2+\|v\|^2.\] 
\end{theorem}

\begin{theorem}[Triangle Inequality] 
If $u$ and $v$ are vectors in an inner product space $V$ then \[\|u+v\|\leq \|u\|+\|v\|.\] 
\end{theorem}

\begin{theorem}[Parallelogram Law] 
If $u$ and $v$ are vectors in an inner product space $V$ then \[\|u+v\|^2+\|u-v\|^2= 2(\|u\|^2+\|v\|^2).\] 
\end{theorem}

\subsection{Orthogonal Projections}
\begin{definition} Given two vectors $\vec{v}$ and $\vec{u}$ in $\mathbb{R}^n$ we define
\[\hat{v}=\left(\dfrac{\vec{v}\cdot \vec{u}}{\vec{u}\cdot \vec{u}}\right)\vec{u}\] to be the \emph{orthogonal projection} of $\vec{v}$ onto $\vec{u}$. 
This is sometimes denoted
\[\hat{v}=\text{proj}_{\vec{u}}(\vec{v}).\]
\begin{center}
\includegraphics[scale=.75]{vector_projection.eps}
\end{center}
\end{definition}
In the above picture the green vector $\vec{v}$ is being projected onto the blue vector $\vec{u}$ and the resulting red vector is the projection. Notice that the projection points in the same direction as $\vec{u}$. \\
\ \\
Moreover, 
\[\vec{z}=(\vec{v}-\hat{v})\bot\, \vec{u}\]
and the length of $\vec{z}=\vec{v}-\hat{v}$ is represented by the dashed line.
\begin{definition}
Let $\vec{v}$ and $\vec{u}$ be two vectors let $\hat{v}$ be the projection of $\vec{v}$ onto $\vec{u}$. If we define $\vec{z}=\vec{v}-\hat{v}$ then
\[\vec{v}=\hat{v}+\vec{z}\] where $\hat{v}$ is parallel to $\vec{u}$ and $\vec{z}$ is perpendicular to $\vec{u}$. We call $\hat{v}$ the component of $\vec{v}$ parallel to $\vec{u}$ and $\vec{z}$ is the component of $\vec{v}$ perpendicular to $\vec{u}$. 
\end{definition}
\begin{definition} Let $u$ and $v$ be non-zero vectors in an inner product space $V$. The orthogonal decomposition of $u$ with respect to $v$ is defined to be 
\[u=\frac{\ip{u}{v}}{\ \|v\|^2}v+\left(u-\frac{\ip{u}{v}}{\|v\|^2}v\right)\]
where $\frac{\ip{u}{v}}{\ \|v\|^2}v$ is parallel to $v$ and $\left(u-\frac{\ip{u}{v}}{\ \|v\|^2}v\right)$ is perpendicular to $v$. 
\end{definition}
\begin{theorem}[Cauchy-Schwarz Inequality]If $u$ and $v$ are vectors in an inner product space $V$ then 
\[\left|\ip{u}{v}\right|^2\leq \|u\|\ \|v\|.\]
\end{theorem}
\end{document}








